\section{A feasible strict direction}
\label{q3:strict:section}

The separation argument needs one variation that makes every active
dominant-energy inequality and every boundary expansion strict.
We construct that direction on the same exterior.  Throughout this
section $(N,g,K)$ is a smooth one-ended exterior in
Definition~\ref{q3:intro:exterior-class}, with compact boundary $B$,
normal $\nu$ pointing into $N$, and the decay, integrability and
dominant-energy hypotheses of Proposition~\ref{q3:intro:exterior-inequality}.
No assumption on its ADM causal character is needed.
The construction also gives the exact charge change and controls every
enclosing cut, including cuts touching the boundary.


Conformal strictification by a spacetime Poisson equation was developed
by Jaracz~\cite{Jaracz2025StrictDEC}. Here a regularized drift equation
with nonzero inner Neumann data also gives strict trapping, an exact
energy slope, unchanged momentum, and comparison of all enclosing cuts.
These quantitative conclusions are established below.

\begin{lemma}[Conformal constraint formulas]
\label{q3:reduction:conformal-formulas}
For a smooth function $f$, set
\[
 g_f=e^{4f}g,\qquad K_f=e^{2f}K.
\]
The corresponding constraint densities satisfy
\begin{align}
 16\pi e^{4f}\mu_f
   &=16\pi\mu-8\Delta_g f-8|df|_g^2,
       \label{q3:reduction:conformal-energy}\\
 8\pi J_f
   &=e^{-2f}\bigl(8\pi J+4K(\nabla f,\cdot)\bigr).
       \label{q3:reduction:conformal-momentum}
\end{align}
Here the second identity is an identity of covectors.  Consequently,
\begin{align}
 16\pi e^{4f}(\mu_f-|J_f|_{g_f})
 \ge{}&16\pi(\mu-|J|_g)\notag\\
 &+8\bigl(-\Delta_g f-|df|_g^2-|K|_g|df|_g\bigr).
       \label{q3:reduction:conformal-dec}
\end{align}
On the fixed boundary,
\begin{equation}
 \theta_{+,f}=e^{-2f}(\theta_++4\partial_\nu f).
       \label{q3:reduction:conformal-expansion}
\end{equation}
Equivalently, if $u=e^f>0$, the last summand in
Equation~\eqref{q3:reduction:conformal-dec} is
$8u^{-1}(-\Delta_g u-|K|_g|du|_g)$.
\end{lemma}

\begin{proof}
The scalar-curvature transformation in dimension three is
$e^{4f}\Scal_{g_f}=\Scal_g-8\Delta_g f-8|df|^2$.
Both quadratic terms in $K$ scale by $e^{-4f}$, which proves
Equation~\eqref{q3:reduction:conformal-energy}.
Put $\pi=K-(\tr_gK)g$.  Then
$\pi_f=e^{2f}\pi$.  For a covariant symmetric tensor $T$, the connection
difference for $\widetilde g=\Omega^2g$ gives, in dimension three,
\[
 \widetilde\nabla^{,i}(\Omega T)_{ij}
 =\Omega^{-1}\bigl(\nabla^iT_{ij}
          +2T_{ij}\nabla^i\log\Omega
          -(\tr_gT)\partial_j\log\Omega\bigr).
\]
Take $\Omega=e^{2f}$ and use $\tr_g\pi=-2\tr_gK$.
The terms containing $\tr_gK$ cancel, leaving
Equation~\eqref{q3:reduction:conformal-momentum}.  Taking its norm introduces
another factor $e^{-2f}$.  The triangle inequality, with the factor two
between the constraint normalizations $16\pi$ and $8\pi$, yields
Equation~\eqref{q3:reduction:conformal-dec}.
Finally, $\nu_f=e^{-2f}\nu$,
$H_f=e^{-2f}(H+4\partial_\nu f)$, and
$\tr_{B,g_f}K_f=e^{-2f}\tr_{B,g}K$.  This proves
Equation~\eqref{q3:reduction:conformal-expansion}.  The last assertion follows
from $\Delta\log u+|d\log u|^2=u^{-1}\Delta u$.
\end{proof}

\begin{lemma}[Strict direction under the original decay]
\label{q3:reduction:strict-direction}
Choose
\[
 0<\delta<\beta<\min(q,1).
\]
There are smooth positive functions $w$ and $\phi$ on $N$, with
$w=r^{-3-\delta}$ on the end, such that
\begin{align}
 &\partial_\nu\phi=-1\quad\hbox{on }B,
       &\phi&=O_2(r^{-1}),\label{q3:reduction:strict-boundary}\\
 &-\Delta_g\phi-|K|_g|d\phi|_g\ge w,
       &\frac{-\Delta_g\phi}{w}&\longrightarrow1.
       \label{q3:reduction:strict-source}
\end{align}
The Euclidean normal flux
\[
 L_\phi=\lim_{r\to\infty}
          \int_{S_r}\partial_r\phi\,\dd A_\delta
\]
is finite and negative.  For every finite $s\ge0$, set
\begin{equation}
 \label{q3:reduction:linear-family}
 u_s=1+s\phi,\qquad g_s=u_s^4g,\qquad K_s=u_s^2K.
\end{equation}
These data are complete up to $B$ and satisfy DEC, strictly for $s>0$.
If $\theta_+\le0$ on $B$, their boundary expansion remains nonpositive
and is strictly negative for $s>0$.  Their decay exponent can be taken to be
$\min(q,1)>1/2$, their constraint densities are integrable, and
\begin{equation}
 \label{q3:reduction:strict-charges}
 E_s=E+sA_\phi,\qquad P_s=P,
 \qquad A_\phi=-\frac{L_\phi}{2\pi}>0.
\end{equation}
Their full enclosing infima satisfy
\begin{equation}
 \label{q3:reduction:strict-areas}
 a_g(B)\le a_{g_s}(B)
       \le (1+s\|\phi\|_\infty)^4a_g(B).
\end{equation}
In particular $E_s\to E$ and $a_{g_s}(B)\to a_g(B)$ as $s\downarrow0$.
No sign assumption on the ADM vector is needed.

If $K$ is compactly supported and the metric has symbol estimates of all
orders with $g-\delta=O(r^{-1})$, the function can instead be chosen with
$\phi=O_k(r^{-1})$ for every $k$, while retaining
Equations~\eqref{q3:reduction:strict-boundary} and
\eqref{q3:reduction:strict-source}.
\end{lemma}

\begin{proof}
Choose smooth $b\ge|K|_g$, positive everywhere and equal to
$b_0r^{-1-\beta}$ sufficiently far out, where $b_0>0$ is chosen large
enough.  This is possible because $\beta<q$.  Choose a smooth positive
regularizer $\zeta$, equal to $r^{-2}$ on the end, and extend $w$ smoothly
and positively to $N$.  We solve
\begin{equation}
 \label{q3:reduction:drift-poisson}
 -\Delta_g\phi=b\sqrt{|d\phi|_g^2+\zeta^2}+w,
 \qquad \partial_\nu\phi=-1,
 \qquad \phi\longrightarrow0.
\end{equation}
The regularizer makes this equation smooth even at critical points.

First truncate at a large coordinate sphere $S_R$ and prescribe
$\phi=0$ there.  The Neumann and Dirichlet conditions occur on disjoint
boundary components.  For $0\le t\le1$, replace $b$ in
Equation~\eqref{q3:reduction:drift-poisson} by $tb$.  At $t=0$, the mixed
Laplacian is invertible: its Dirichlet face gives the Poincare inequality,
and the weak solution obtained from its coercive quadratic form is smooth
by boundary regularity.  At any solution, the linearization is a
Laplacian with a smooth drift of norm at most $b$ and with the homogeneous
mixed boundary conditions.  The strong maximum principle and the Hopf
boundary principle make its kernel zero.  It is a Fredholm perturbation
of the mixed Laplacian of index zero, and is therefore invertible.
The boundary estimates used here are the ordinary elliptic estimates
for Dirichlet and normal Neumann conditions: positive definiteness of
$g$ verifies their principal complementing condition, and the disjoint
boundary components can be covered separately
\cite{AgmonDouglisNirenberg1959}.  The maximum principles and local
elliptic estimates below are used in their usual uniformly elliptic
forms \cite{GilbargTrudinger2001}.

To close continuation at $t=1$, first work on a fixed truncation.
The elliptic $W^{2,p}$ estimate
and gradient interpolation give
\[
 \|\phi\|_{W^{2,p}}
 \le C_p\bigl(1+\|\phi\|_{L^p}+\|d\phi\|_{L^p}\bigr)
 \le \tfrac12\|\phi\|_{W^{2,p}}
       +C'_p(1+\|\phi\|_{L^p}).
\]
The constants include the fixed boundary data and are uniform in $t$.
If the supremum norms of solutions on this truncation were unbounded,
divide them by their supremum norms.  The preceding estimate with
$p>3$, compactness, and the equation give a nonzero limit $v$ with
homogeneous mixed data satisfying
\[
 -\Delta_gv=t_*b|dv|_g.
\]
This can be written $\Delta_gv+A\cdot dv=0$ with a bounded measurable
drift: take $A=t_*b\nabla v/|dv|$ off the critical set and zero on it.
The strong and Hopf maximum principles exclude a nonzero solution with
the homogeneous mixed conditions.  Thus the supremum norms are bounded.
The $W^{2,p}$ estimate, Schauder estimates, and the smooth equation now
give the bounds needed for closedness of continuation.  Existence on
each truncation follows.

Every such solution is nonnegative.  A negative minimum cannot be
interior because the right side of
Equation~\eqref{q3:reduction:drift-poisson} is positive.  At an inner-boundary
minimum, the outward-domain derivative would be negative by the Hopf
principle, whereas the prescribed derivative in that direction is $1$.
The outer value is zero.

It remains to make these bounds independent of $R$.  For fixed sufficiently large
$A$ and then sufficiently large $r_0$, the positive radial function
\[
 v_0(r)=r^{-1}(1-Ar^{-\delta})
\]
satisfies, on $r\ge r_0$,
\begin{align*}
 -\Delta_gv_0-b|dv_0|_g
 &=A\delta(1+\delta)r^{-3-\delta}
       +O(r^{-3-q})+O(r^{-3-\beta})\ge c_0w.
\end{align*}
The constants are fixed independently of the truncation.  Since
$b\zeta=o(w)$, a sufficiently large multiple of $v_0$ is a supersolution
of the regularized equation.  The comparison principle gives
\begin{equation}
 \label{q3:reduction:core-barrier}
 0\le\phi(x)\le C(1+M_R)v_0(r),\qquad
 M_R=\sup_{N\cap\{r\le r_0\}}\phi,
 \quad r_0\le r\le R.
\end{equation}
The notation for the compact core includes all points outside the end
chart.  Were $M_R$ unbounded along expanding truncations, the normalized
solutions $\phi/M_R$ would have locally uniform $W^{2,p}$ bounds, including
at $B$, by the preceding local estimates and
Equation~\eqref{q3:reduction:core-barrier}.  A subsequence would converge
locally in $C^1$ to a nonnegative function with maximum one on the core,
homogeneous inner Neumann data, and a bounded-drift homogeneous equation.
Equation~\eqref{q3:reduction:core-barrier} makes its value tend to zero at
infinity.  It consequently attains a positive maximum, contradicting the
strong or Hopf principle for that homogeneous equation.  This proves the
uniform core bound.  The contradiction uses the homogeneous bounded-drift
equation.

Local compactness and diagonal extraction give a smooth solution of
Equation~\eqref{q3:reduction:drift-poisson}, with $\phi=O(r^{-1})$.
The claimed derivative count can be checked directly on annuli.  Put
\[
 \phi_{\rho}(y)=\rho\phi(\rho y),\qquad
 g_{\rho}(y)=g(\rho y),\qquad 1<|y|<4.
\]
On smaller fixed annuli the equation is
\[
 -\Delta_{g_{\rho}}\phi_{\rho}
 =\rho b(\rho y)
       \sqrt{|d\phi_{\rho}|_{g_{\rho}}^2+
                     (\rho^2\zeta(\rho y))^2}
       +\rho^3w(\rho y).
\]
Its coefficients have uniform $C^{1,1}$ bounds from the two derivatives
in Equation~\eqref{q3:intro:decay}; the scaled drift is
$O(\rho^{-\beta})$ and the scaled source is $O(\rho^{-\delta})$.
The bounded zeroth-order norm, interior $W^{2,p}$ estimates, and gradient
interpolation first give uniform $C^{1,\alpha}$ bounds for some
$\alpha>0$.  The right side is then uniformly $C^\alpha$.  Schauder
estimates give uniform $C^{2,\alpha}$ bounds on smaller annuli.  This
proves $\phi=O_2(r^{-1})$ without using a third derivative of $g$ or a
second derivative of $K$.

The right side of Equation~\eqref{q3:reduction:drift-poisson} is
$w+O(r^{-3-\beta})$.  It is integrable and proves
Equation~\eqref{q3:reduction:strict-source}.  The divergence theorem, with
outward-domain normal $-\nu$ on $B$, gives
\begin{equation}
 \label{q3:reduction:phi-flux}
 \lim_{r\to\infty}\int_{S_r}\partial_{n_g}\phi\,\dd A_g
 =-|B|_g-\int_N
       \bigl(b\sqrt{|d\phi|_g^2+\zeta^2}+w\bigr)\,\dd V_g.
\end{equation}
The Euclidean flux differs by $O(r^{-q})$ and has the same finite limit.

For every fixed finite $s\ge0$, the factor $u_s$ in
Equation~\eqref{q3:reduction:linear-family} is bounded and at least one,
so the transformed data are complete up to $B$.
The $u$-form of Lemma~\ref{q3:reduction:conformal-formulas} gives
\begin{align*}
 16\pi u_s^4(\mu_s-|J_s|_{g_s})
 &\ge16\pi(\mu-|J|_g)
        +\frac{8s}{u_s}(-\Delta_g\phi-|K|_g|d\phi|_g)
 \ge \frac{8s}{u_s}w,\\
 \theta_{+,s}
 &=u_s^{-2}\left(\theta_+-\frac{4s}{u_s}\right).
\end{align*}
This proves the asserted signs for all finite $s$, without a smallness
restriction.

Since $u_s-1=O_2(r^{-1})$, the transformed decay exponent is
$\min(q,1)>1/2$.  The exact energy transformation is
\[
 16\pi u_s^4\mu_s=16\pi\mu-8s u_s^{-1}\Delta_g\phi.
\]
Here $\Delta_g\phi=O(r^{-3-\delta})$, and the additional current term
is bounded by a fixed-$s$ constant times
$|K|_g|d\phi|_g=O(r^{-3-q})$.  Both are integrable in dimension three.
The bounded positive conformal factor also preserves integrability in
the transformed volume measure.

The leading metric change is $4s\phi\delta$, whose ADM energy flux
is $-8s\int_{S_r}\partial_r\phi\,\dd A_\delta$ before division by
$16\pi$.  All remaining metric flux terms are
$O_s(r^{-q})+O_s(r^{-1})$.  Writing $\pi=K-(\tr_gK)g$, one has
$\pi_s=u_s^2\pi$, so the change in its Euclidean momentum flux is
$O_s(r^2r^{-1}r^{-1-q})=O_s(r^{-q})$.
This proves Equation~\eqref{q3:reduction:strict-charges};
Equation~\eqref{q3:reduction:phi-flux} gives $A_\phi>0$.
No $1/r$ coefficient expansion for $\phi$ is needed.
Finally, the cut class is unchanged and every cut has area
$\int_\Gamma u_s^4\,\dd A_g$.  Bounding $u_s$ between $1$ and
$1+s\|\phi\|_\infty$ and taking infima proves
Equation~\eqref{q3:reduction:strict-areas}, including coincident and
disconnected cuts.

For the final assertion, take $b$ nonnegative, compactly supported, and
at least $|K|$; the proof is unchanged.  The end equation is then simply
$-\Delta_g\phi=w$.  Once the two-derivative estimate has been established,
successive annular elliptic estimates use the assumed symbol bounds for
$g$ and $w$ to give $\phi=O_k(r^{-1})$ for every $k$.
\end{proof}

